\section{Section de charge et relations de quantification électrique}
\label{iii:sec:carga}

La réponse constitutive établit une relation entre les deux feuillets et une échelle relative de propagation. La section de charge s'obtient en composant cette réponse avec la constante de structure fine et l'action, toutes deux construites dans les chapitres précédents. L'ordre a une conséquence précise : la charge ne détermine pas rétrospectivement le registre dodécaphasique. C'est le registre, avec l'action et la réponse du vide, qui détermine la section de charge dans la droite correspondante.

\subsection{Détermination positive de la section de charge}

Considérons les sections positives $\hbar_a$, avec
$a\in\{\mathrm{pre},\mathrm{ret}\}$, et $h_a=2\pi\hbar_a$.
L’impédance $Z_{\rm vac}$ est déjà construite au moyen de
\eqref{iii:eq:dimensions}. Le quotient $h_a/Z_{\rm vac}$ appartient à
la droite quadratique de charge. Cette compatibilité dimensionnelle permet de
formuler la définition suivante sans choisir encore de coordinatisation SI.

\begin{definition}[Section positive de charge]
La section de charge correspondant à l’orientation d’action $a$ est
\begin{equation}
 e_a=\left(\frac{2\alpha h_a}{Z_{\rm vac}}\right)^{1/2}.
 \label{iii:eq:charge}
\end{equation}
La racine est prise dans la demi-droite positive de charge. Sa section conjuguée
est $-e_a$ ; le signe de cette dernière n’est pas identifié à l’indice de
retour de l’action.
\end{definition}

\begin{proposition}[Existence, unicité et compatibilité du couplage]
L’équation $Z_{\rm vac}e_a^2=2\alpha h_a$ détermine une unique section
positive. Dans la réalisation constitutive, on a aussi
\begin{equation}
 \alpha=\frac{Z_{\rm vac}e_a^2}{2h_a}
 =\frac{e_a^2}{4\pi\varepsilon_{\rm vac}\hbar_a c_{\rm vac}}.
 \label{iii:eq:coupling}
\end{equation}
\end{proposition}
\begin{proof}
La positivité de $\alpha$, $h_a$ et $Z_{\rm vac}$ garantit une unique racine positive. Les identités constitutives donnent $\varepsilon_{\rm vac}c_{\rm vac}=Z_{\rm vac}^{-1}$. Substituer cette égalité et $h_a=2\pi\hbar_a$ dans la première expression de \eqref{iii:eq:coupling} produit la seconde.
\end{proof}

La dernière expression est la forme habituelle du couplage électromagnétique
dans une coordinatisation physique. Elle apparaît ici comme identité ultérieure des sections
construites. La preuve ne recalcule pas $\alpha$ à partir d’une charge
mesurée : elle établit la cohérence algébrique du résultat précédent avec la
section qui vient d’être définie. L’identification physique exige en outre
de maintenir la coordinatisation des unités et le régime du couplage employés dans la
comparaison. Une identité entre grandeurs et une comparaison expérimentale
ont des fonctions différentes dans cette reconnaissance.

\subsection{Résistance, conductance, flux et fréquence}

Les quatre grandeurs suivantes utilisent la même section $e_a$. Leur
dépendance peut être entièrement résolue avant l’évaluation de la moindre décimale.

\begin{theorem}[Composition des quanta électriques]
Les sections
\begin{align}
 R_K&=\frac{h_a}{e_a^2}=\frac{Z_{\rm vac}}{2\alpha},
 &G_0&=\frac{2e_a^2}{h_a}=\frac{4\alpha}{Z_{\rm vac}},
 \label{iii:eq:rk-g0}\\
 \Phi_{0,a}&=\frac{h_a}{2e_a}
 =\sqrt{\frac{h_aZ_{\rm vac}}{8\alpha}},
 &K_{J,a}&=\frac{2e_a}{h_a}=\Phi_{0,a}^{-1}
 \label{iii:eq:flux-josephson}
\end{align}
satisfont exactement
\begin{equation}
 R_KG_0=2,\qquad G_0Z_{\rm vac}=4\alpha,\qquad
 K_{J,a}\Phi_{0,a}=1,\qquad
 K_{J,a}^{2}R_K=\frac4{h_a}.
 \label{iii:eq:quanta-identities}
\end{equation}
\end{theorem}
\begin{proof}
La substitution de $e_a^2=2\alpha h_a/Z_{\rm vac}$ dans les deux premières
définitions annule $h_a$ et donne \eqref{iii:eq:rk-g0}. En élevant au
carré la section positive $h_a/(2e_a)$, on obtient
$h_aZ_{\rm vac}/(8\alpha)$ ; la positivité fixe la racine dans
\eqref{iii:eq:flux-josephson}. Son inverse est $2e_a/h_a$.
Les trois premiers produits de \eqref{iii:eq:quanta-identities} se
réduisent, respectivement, à $2$, $4\alpha$ et $1$. Enfin,
$(4e_a^2/h_a^2)(h_a/e_a^2)=4/h_a$.
\end{proof}

La réalisation physique reconnaît dans ces sections la résistance de von Klitzing, le quantum de conductance, le quantum de flux et la constante de Josephson. Le facteur $2$ de la conductance appartient à la normalisation de la grandeur définie ici ; son identification à une multiplicité de canaux d'un système matériel exige de spécifier cette représentation.
Les quatre équations ne constituent pas quatre générateurs indépendants : elles partagent $\alpha$, l'action et la réponse du vide. C'est précisément cette dépendance qui produit les identités croisées et permet de vérifier une évaluation sans altérer les opérations qui l'ont engendrée.

\subsection{Retour de l’action et transport des sections}

Le quotient $R_{\rm act}=h_{\rm pre}/h_{\rm ret}$, obtenu avant le
couple angulaire, conserve la différence entre les deux sections d’action.
Il ne change pas les assignations de feuillet de la réponse constitutive. Ce sont deux
opérations distinctes : retour de l’action et échange des canaux.

\begin{proposition}[Covariance sous le retour]
Avec $\alpha$ et $Z_{\rm vac}$ communs, on a
\begin{align}
 \frac{e_{\rm pre}}{e_{\rm ret}}&=R_{\rm act}^{1/2},
 &R_K^{\rm pre}&=R_K^{\rm ret},
 &G_0^{\rm pre}&=G_0^{\rm ret},\label{iii:eq:charge-return}\\
 \frac{\Phi_{0,\rm pre}}{\Phi_{0,\rm ret}}&=R_{\rm act}^{1/2},
 &\frac{K_{J,\rm pre}}{K_{J,\rm ret}}&=R_{\rm act}^{-1/2}.
 \label{iii:eq:flux-return}
\end{align}
\end{proposition}
\begin{proof}
La première identité découle de la division des deux versions de
\eqref{iii:eq:charge} et de la prise de la racine positive. Dans
\eqref{iii:eq:rk-g0}, l’action s’est déjà annulée. Les deux autres
identités découlent des puissances $h_a^{1/2}$ et $h_a^{-1/2}$ dans
\eqref{iii:eq:flux-josephson}.
\end{proof}

La résistance et la conductance conservent le quotient constitutif, tandis que le flux et la fréquence retiennent également l'orientation de l'action. Cette séparation donne une lecture structurelle de leurs dépendances : certaines grandeurs sont des invariants du quotient des feuillets et d'autres mesurent en outre la section d'action utilisée. La comparaison métrologique doit déclarer quelle section elle évalue ; la proximité numérique des deux actions n'autorise pas à les identifier.
